\documentclass[10pt,a4paper]{article}
\usepackage[utf8]{inputenc}
\usepackage[indonesian]{babel}
\usepackage{amsmath,amssymb,amsfonts}
\usepackage{geometry}
\usepackage{multicol}
\usepackage{enumitem}
\usepackage{fourier}

\geometry{top=1.5cm, bottom=1.5cm, left=1.5cm, right=1.5cm}

\pagestyle{empty}

\begin{document}

\begin{center}
  \textbf{\Large Latihan Soal PPKN}
\end{center}
\vspace{0.3cm}

\begin{multicols}{2}
  \setlist[enumerate,1]{leftmargin=*,label=\arabic*.}
  \setlist[enumerate,2]{leftmargin=*,label=\Alph*.}

  \begin{enumerate}
  \item Kehidupan yang rukun dan damai dengan sikap dan perbuatan yang menunjukkan toleransi dalam keberagaman didambakan oleh setiap orang. Hal tersebut untuk mewujudkan\dots
    \begin{enumerate}
    \item Kemakmuran
    \item Kemerdekaan
    \item Keharmonisan
    \item Gotong royong
    \item \textbf{Persatuan dan kesatuan}
    \end{enumerate}

  \item Pada awal persidangan, Ketua BPUPK mengajukan pertanyaan mengenai, "Apa yang akan menjadi dasar negara Indonesia merdeka?" Ketua BPUPK adalah\dots
    \begin{enumerate}
    \item Supomo
    \item Ir. Sukarno
    \item Ki Bagoes Hadikoesoemo
    \item Drs. Mohammad Hatta
    \item \textbf{KRT Radjiman Wedyodiningrat}
    \end{enumerate}

  \item Rumusan dasar negara Pancasila yang digali oleh Ir. Sukarno berasal dari\dots
    \begin{enumerate}
    \item Kehidupan suku-suku yang berkembang di Nusantara
    \item Sintesis ideologi-ideologi besar di dunia pada waktu itu
    \item Tata kehidupan masyarakat pada masa kerajaan di Nusantara
    \item Budaya asing yang memperkuat budaya dan nilai-nilai dalam masyarakat Indonesia
    \item \textbf{Budaya, adat istiadat, agama, dan nilai-nilai yang hidup dalam masyarakat Indonesia}
    \end{enumerate}

  \item Persidangan BPUPK berada di bawah tekanan Pemerintahan Militer Jepang. Meskipun demikian, para tokoh anggota BPUPK tetap melaksanakan tugas untuk kepentingan masa depan bangsa yang merdeka. Nilai-nilai luhur yang dapat kalian teladani dari pernyataan tersebut adalah\dots
    \begin{enumerate}
    \item Keberanian
    \item \textbf{Kerja keras}
    \item Rendah hati
    \item Rela berkorban
    \item Menghargai pendapat
    \end{enumerate}

  \item Upaya dalam mengelola perbedaan pendapat yang tajam terhadap Piagam Jakarta dilakukan oleh para tokoh seperti Muhammad Hatta dan Sukarno dengan melaukan pertemuan di luar forum sidang dengan berbagai tokoh dan kalangan. Upaya tersebut dikenal dengan cara\dots
    \begin{enumerate}
    \item Seminar
    \item Diskusi ilmiah
    \item \textbf{Rapat di luar sidang}
    \item Komunikasi dan lobi
    \item Pendekatan persuasif
    \end{enumerate}

  \item Piagam Jakarta mendapat masukan dan tanggapan dari para tokoh anggota BPUPK dan PPKI. Berbagai perubahan dilakukan sehingga berhasil dirumuskan dasar negara Pancasila dan ditetapkan oleh PPKI pada tanggal 18 Agustus 1945. Kalian dapat meneladani nilai-nilai luhur dalam proses tersebut, yaitu\dots
    \begin{enumerate}
    \item \textbf{Menerima dan melaksanakan hasil keputusan bersama}
    \item Perbedaan pendapat berlangsung meriah dan menarik untuk dikaji
    \item Untuk kepentingan bersama maka perlu terjadi perbedaan pendapat
    \item Mempermasalahkan perubahan karena tidak sesuai kesepakatan awal
    \item Banyak usulan maka sulit untuk dipersatukan dalam sebuah kesepakatan bersama
    \end{enumerate}

  \item Rumusan Pancasila pada Piagam Jakarta mengalami perubahan, yaitu pencoretan atau penghapusan pada tujuh kata, yaitu dengan kewajiban menjalankan syariat Islam bagi pemeluk-pemeluknya. Sebelumnya, tujuh kata tersebut melekat pada sila\dots dalam Piagam Jakarta.
    \begin{enumerate}
    \item \textbf{Ketuhanan}
    \item Kemanusiaan
    \item Persatuan
    \item Kerakyatan
    \item Keadilan sosial
    \end{enumerate}

  \item Rumusan dasar negara sebagaimana tercantum dalam UUD NRI Tahun 1945 hasil penetapan PPKI tanggal 18 Agustus 1945 mengalami perubahan naskah pada masa RIS dan UUDS 1950. Namun, dapat dikatakan bahwa\dots
    \begin{enumerate}
    \item Nilai-nilainya kembali kepada Piagam Jakarta 22 Juni 1945
    \item \textbf{Nilai-nilai dasar dipengaruhi oleh bentuk negara yang berlaku}
    \item Kandungan nilai (prinsip-prinsip pokok) setiap sila mengalami perubahan
    \item Nilai-nilainya dipengaruhi oleh konsep demokrasi yang berlaku, yaitu liberal
    \item Kandungan nilai (prinsip-prinsip pokok) setiap sila secara substantif tidak berubah
    \end{enumerate}

  \item Rumusan dasar negara Pancasila yang berlaku sekarang adalah rumusan yang ditetapkan berdasarkan Dekrit Presiden 5 Juli 1959 dan Lembaran Negara RI Nomor 75 Tahun 1959. Meskipun demikian, ada perubahan teks naskah yang berbeda dengan rumusan PPKI tanggal 18 Agustus 1945, yaitu pada kalimat\dots
    \begin{enumerate}
    \item \textbf{Ketuhanan Yang Maha Esa}
    \item Kemanusiaan
    \item Persatuan Indonesia
    \item Permusyawaratan/Perwakilan
    \item Keadilan sosial
    \end{enumerate}

  \item Perhatikan pernyataan berikut:
    \begin{enumerate}[label=(\arabic*)]
    \item Bersendikan nilai-nilai agama.
    \item Mengandung unsur-unsur budaya dan adat istiadat.
    \item Menjadi dasar alasan filosofis berdirinya suatu negara.
    \item Setiap produk hukum di Indonesia harus berdasarkan Pancasila.
    \end{enumerate}
    Berdasarkan pernyataan tersebut, Pancasila merupakan dasar filsafat negara (\textit{philosophische grondslag}) karena alasan yang ditunjukkan pada nomor\dots
    \begin{enumerate}
    \item (1) dan (2)
    \item (1) dan (3)
    \item (2) dan (3)
    \item (2) dan (4)
    \item \textbf{(3) dan (4)}
    \end{enumerate}

  \item Transformasi dilakukan untuk perubahan ke arah yang lebih baik lagi. Bangsa Indonesia yang berkepribadian dalam kebudayaan merupakan tujuan dari transformasi\dots
    \begin{enumerate}
    \item Politik
    \item Material
    \item Ekonomi
    \item \textbf{Mental karakter}
    \item Pertahanan dan keamanan
    \end{enumerate}

  \item Hubungannya dengan pemenuhan cita-cita (visi) dan tujuan Negara Indonesia dari Pancasila sebagai ideologi negara maka Mohammad Hatta memerasnya atau menyimpulkan menjadi\dots
    \begin{enumerate}
    \item \textbf{Persatuan}
    \item Perdamaian
    \item Kebahagiaan
    \item Kemerdekaan
    \item Kesejahteraan
    \end{enumerate}

  \item Perwujudan Pancasila sebagai dasar filsafat negara memerlukan prasyarat kecerdasan. Oleh karena itu, harus ada upaya yang sungguh-sungguh untuk mewujudkan tujuan negara, terutama\dots
    \begin{enumerate}
    \item Memajukan kesejahteraan umum
    \item \textbf{Mencerdaskan kehidupan bangsa}
    \item Ikut melaksanakan ketertiban dunia
    \item Mewujudkan masyarakat adil dan makmur
    \item Melindungi segenap bangsa dan seluruh tumpah darah Indonesia
    \end{enumerate}

  \item Warga negara dituntut bukan saja memiliki kecerdasan dan karakter personal, namun diutamakan kecerdasan dan karakter kewargaan dalam menyikapi keberagaman masyarakat, yaitu\dots
    \begin{enumerate}
    \item Kompeten mengemban tugas kewargaan
    \item \textbf{Mampu mencari titik temu dalam perbedaan}
    \item Memahami hak dan kewajiban sebagai warga negara
    \item Mampu menempatkan keunggulan pribadi dalam harmoni kemajuan bersama
    \item Mampu memenuhi panggilan keterlibatan dalam urusan publik secara sukacita
    \end{enumerate}

  \item Contoh implementasi nilai kemanusiaan dalam praktik kehidupan di sekolah adalah\dots
    \begin{enumerate}
    \item Pemilihan Ketua OSIS
    \item Pelaksanaan piket kelas
    \item Pemilihan pelajar berprestasi
    \item Kompetisi sains dan olahraga
    \item \textbf{Donasi bagi korban bencana alam}
    \end{enumerate}

  \item Bintang penuntun atau \textit{leitstar} bagi upaya mengisi kemerdekaan dengan pembangunan nasional, yaitu\dots
    \begin{enumerate}
    \item \textbf{Pancasila}
    \item Bahasa Indonesia
    \item Bendera Merah Putih
    \item Bhinneka Tunggal Ika
    \item UUD NRI Tahun 1945
    \end{enumerate}

  \item Berikut yang merupakan upaya pembudayaan nilai-nilai Pancasila melalui peringatan jasa para pahlawan adalah\dots
    \begin{enumerate}
    \item Memberikan keteladanan atau contoh baik
    \item Mengetahui daftar nama pahlawan dan jasa-jasanya
    \item Mengingat perjuangan para pahlawan agar terinspirasi
    \item \textbf{Meneruskan semangat terus berjuang melawan penjajah}
    \item Nostalgia bagaimana cara melawan penjajahan pada masa lalu
    \end{enumerate}

  \item Nilai Pancasila, yaitu Keadilan Sosial diterapkan dalam kebijakan pendidikan dengan cara\dots
    \begin{enumerate}
    \item Membatasi akses pembangunan pendidikan untuk kelompok tertentu
    \item Meniadakan pendidikan untuk kelompok minoritas dan daerah tertinggal
    \item Mengabaikan hak-hak sosial dan kearifan lokal dalam konteks pendidikan
    \item \textbf{Memastikan distribusi pendidikan yang adil bagi segenap lapisan masyarakat}
    \item Menekankan pada pendidikan tertentu dan meniadakan pendidikan untuk semua
    \end{enumerate}

  \item Nilai Pancasila, yaitu gotong royong dapat tercermin dalam penanganan bencana alam yang terjadi dalam masyarakat dengan cara\dots
    \begin{enumerate}
    \item Mendorong persaingan dalam memperoleh bantuan
    \item Menciptakan ketidaksetaraan dalam distribusi bantuan
    \item Menekankan pada tanggung jawab masing-masing individu
    \item Menunggu dan mengandalkan bantuan dari pemerintah untuk masyarakat
    \item \textbf{Membangun kerja sama dan saling membantu di antara warga masyarakat}
    \end{enumerate}

  \item Untuk mewujudkan kemaslahatan dan kebahagiaan hidup bersama maka ditempuh revitalisasi pembangunan melalui jalur kerukunan hidup antarwarga masyarakat dan bangsa yang mengarah pada terwujudnya Negara Indonesia\dots
    \begin{enumerate}
    \item \textbf{Terpuji}
    \item Bersatu
    \item Berbagi sejahtera
    \item Cerdas kewargaan
    \item Tertata-terlembaga
    \end{enumerate}
  \end{enumerate}

\end{multicols}

\end{document}
